\documentclass[12pt]{article}
\usepackage[T1]{fontenc}
\usepackage[utf8]{inputenc}
\usepackage{lmodern}
\usepackage{textcomp}
\usepackage{enumitem}
\usepackage{geometry}
\geometry{margin=1in}
\begin{document}
\begin{center}
Answer Key - Physics - Undergraduate Level Assessment 16 \
\small (Answers are ordered exactly as in the question paper.)
\end{center}

\section*{Multiple Choice}
\begin{enumerate}[label=\arabic*.]
\item \textbf{Answer:} D
\\ \textit{Options:} A) 0.1 GeV/$c^2$; B) 0.2 GeV/$c^2$; C) 0.5 GeV/$c^2$; D) 1.0 GeV/$c^2$
\\ \textit{Note:} The correct answer is 1.0 GeV/c² because, using the relativistic energy-momentum relation E² = (pc)² + (m₀c²)², we can calculate the rest mass m₀ from the given total energy and momentum, resulting in approximately 1.0 GeV/c².
\item \textbf{Answer:} C
\\ \textit{Options:} A) electrons have many more degrees of freedom than atoms do; B) the electrons and the lattice are not in thermal equilibrium; C) the electrons form a degenerate Fermi gas; D) electrons in metals are highly relativistic
\\ \textit{Note:} The mean kinetic energy of conduction electrons in metals is higher than kT because they occupy energy states up to the Fermi level in a degenerate Fermi gas, resulting in a significant distribution of high-energy electrons even at low temperatures.
\item \textbf{Answer:} C
\\ \textit{Options:} A) 1.2 x $10^9$ m/s; B) 3.0 x $10^8$ m/s; C) 1.5 x $10^8$ m/s; D) 1.0 x $10^8$ m/s
\\ \textit{Note:} The speed of light in a medium is given by the equation c/n, where c is the speed of light in a vacuum (approximately 3 x $10^8$ m/s) and n is the refractive index, which for a dielectric material is equal to the square root of its dielectric constant; thus, with a dielectric constant of 4.0, the refractive index is 2.0, resulting in a speed of light of 1.5 x $10^8$ m/s (3 x $10^8$ m/s divided by 2).
\item \textbf{Answer:} A
\\ \textit{Options:} A) 4; B) 5; C) 6; D) 20
\\ \textit{Note:} The angular magnification of a refracting telescope is given by the ratio of the focal lengths of the objective lens to the eyepiece lens, and with an objective lens focal length of 80 cm (100 cm - 20 cm for the eyepiece), the calculation yields a magnification of 4 (80 cm / 20 cm).
\item \textbf{Answer:} D
\\ \textit{Options:} A) independent of M; B) proportional to $m^(1$/2); C) linear in R; D) proportional to $R^(3$/2)
\\ \textit{Note:} The correct answer is based on Kepler's third law of planetary motion, which states that the square of the orbital period (T) of a satellite is directly proportional to the cube of the semi-major axis (R) of its orbit, leading to the proportionality T ∝ $R^(3$/2).
\end{enumerate}

\section*{True / False}
\begin{enumerate}[label=\arabic*.]
\setcounter{enumi}{5}
\item \textbf{Answer:} False
\\ \textit{Options:} A) True; B) False
\\ \textit{Note:} The statement is false because a reversible thermodynamic process can occur at varying temperatures, depending on the specific conditions of the system and the nature of the process, such as isothermal, adiabatic, or isochoric processes.
\item \textbf{Answer:} False
\\ \textit{Options:} A) True; B) False
\\ \textit{Note:} The statement is false because an observer moving at 0.50 c would measure the length of the rod to be contracted according to the length contraction formula, which would yield a contraction factor of approximately 0.866, resulting in a measured length of about 0.866 m, not 0.80 m.
\item \textbf{Answer:} False
\\ \textit{Options:} A) True; B) False
\\ \textit{Note:} The statement is false because the Sun's energy is primarily produced by a series of nuclear fusion reactions in which four hydrogen nuclei (protons) ultimately fuse to form one helium nucleus, rather than just two hydrogen atoms.
\item \textbf{Answer:} False
\\ \textit{Options:} A) True; B) False
\\ \textit{Note:} The statement is false because the electric displacement current is related to the rate of change of the electric field, not the magnetic flux, as described by Maxwell's equations.
\item \textbf{Answer:} False
\\ \textit{Options:} A) True; B) False
\\ \textit{Note:} The statement is false because the quantum efficiency of 0.1 indicates that, on average, the detector will detect 10\% of the photons sent, but the actual number of detected photons can vary due to the probabilistic nature of quantum processes.
\end{enumerate}

\section*{Long Form Response}
\begin{enumerate}[label=\arabic*.]
\setcounter{enumi}{10}
\item \textbf{Answer:} The decrease in the universe's temperature from 12 K to 3 K has significant implications for the spatial distribution of galaxies. As the universe cools, the wavelengths of light stretch due to the expansion of space, which affects how we perceive distances. During the earlier phase at 12 K, galaxies were approximately one-quarter as distant from each other compared to their current positions. This relative proximity suggests that the gravitational interactions between galaxies would have been more pronounced in the earlier universe, potentially influencing their formation and clustering. Consequently, understanding this temperature change helps elucidate the evolution of cosmic structures over time.
\\ \textit{Note:} correct because it identifies how the cooling of the universe affects light wavelengths and galaxy distances, indicating that higher temperatures correspond to closer proximity of galaxies due to gravitational interactions, which are essential for understanding their spatial distribution.
\item \textbf{Answer:} The Hall coefficient is a critical parameter in determining the sign of charge carriers in a doped semiconductor. When a magnetic field is applied perpendicular to the current flow in a semiconductor, charge carriers experience a Lorentz force, causing them to accumulate on one side of the material, resulting in a measurable Hall voltage. The sign of this voltage, relative to the current direction, indicates the type of charge carriers: a positive Hall coefficient signifies the dominance of holes (positive charge carriers), while a negative Hall coefficient indicates the presence of electrons (negative charge carriers). Understanding the sign of the Hall coefficient is essential for characterizing semiconductor materials and influences their electrical properties and performance in electronic devices. Thus, the Hall effect not only provides insights into charge carrier types but also informs the design and application of semiconductor technologies.
\\ \textit{Note:} The answer correctly identifies that the Hall coefficient reveals the type of charge carriers in a semiconductor by analyzing the direction of the induced Hall voltage, which results from the Lorentz force acting on the charge carriers in the presence of a magnetic field.
\end{enumerate}

\vfill
\noindent\textit{End of Answer Key}
\end{document}